\documentclass[12pt]{ca-capstone}
\geometry{margin=25mm}
\usepackage{setspace}
\setstretch{1.1}
\capstonerefs{msu-ent,galbraith,knuth-taocp2}
\begin{document}
\capstonetitle{Алгоритмы Евклида, порядки и первообразные корни}
\begin{center}\textbf{Теоретическая справка для преподавателя. Не выдавать студентам.}\end{center}

\section{НОД и коэффициенты Безу}
Бинарный НОД использует $\gcd(2a,2b)=2\gcd(a,b)$, $\gcd(2a,b)=\gcd(a,b)$ для нечётного $b$ и
$\gcd(a,b)=\gcd(a-b,b)$. После общего сдвига вычитания и сдвиги сохраняют НОД. Отдельно
проверяются $(0,0)$, отрицательные аргументы и обмен при $a>b$.

Итеративный расширенный Евклид хранит две пары коэффициентов для текущих остатков. Инвариант:
каждый остаток равен $ax+by$. При нормализации знака меняются знаки обоих коэффициентов. Если
$\gcd(a,n)=1$, то $x=a^{-1}\pmod n$; иначе обратного не существует.

\section{Пакетное обращение}
Для обратимых $a_i$ вычисляются префиксы $P_i=\prod_{j<i}a_j$ и один обратный к полному
произведению $P$. В обратном проходе
\[
 a_i^{-1}=P^{-1}\prod_{j<i}a_j\prod_{j>i}a_j \pmod n,
\]
после чего накопленный обратный умножается на $a_i$. Нули и необратимые элементы заранее
отделяются проверкой НОД и получают \texttt{null}; иначе полное произведение необратимо.
Отдельная проверка НОД не превращает алгоритм в повторное обращение всех элементов.

\section{КТО и Гарнер}
Итеративная КТО: если $x$ решает систему по модулю $M$, добавление $x\equiv r\pmod m$ даёт
\[
 t=(r-x)(M^{-1}\bmod m)\bmod m,\qquad x'=x+Mt,\quad M'=Mm.
\]
Гарнер вычисляет смешанные коэффициенты $c_i$ представления
$x=c_0+c_1m_0+c_2m_0m_1+\cdots$. На шаге $i$ известная часть вычисляется по модулю $m_i$ и
домножается на обратное к $m_0\cdots m_{i-1}$. Взаимная простота модулей обеспечивает все
обратные элементы.

\section{Порядок и дерево произведений}
Если $N=p-1=\prod q^{e_q}$, начинаем с $N$ и для каждого $q$ делим порядок на $q$, пока
$g^{N/q}=1$. Минимальность проверяется всеми простыми делителями исходного $N$, а не пробными
делителями полученного числа.

В дереве листья соответствуют $q^e$. Узлу с элементом $h$ и произведениями поддеревьев $L,R$
левому ребёнку передаётся $h^R$, правому $h^L$. В листе последовательность
$h,h^q,h^{q^2},\ldots$ впервые достигает $1$ при порядке $q^k$. Порядки листьев перемножаются.
Проверка дерева: произведение всех $q^e$ равно $N$; каждая последовательность не длиннее $e+1$.

Элемент $g$ первообразен тогда и только тогда, когда для каждого простого $q\mid p-1$
\[
 g^{(p-1)/q}\not\equiv1\pmod p.
\]
Любой элемент, прошедший критерий, является правильным ответом независимо от выбора стартового
кандидата.

\clearpage
\section{Контрольный список}
Обязательные независимые проверки: линейное тождество Безу; произведение $aa^{-1}$; список
необратимых индексов; все сравнения КТО; $g^r=1$ и минимальность $r$ через Sage; критерий
первообразного корня. Отдельные классы: повторные $q$ в $N$, $g=1$, порядок $2$, большой
простой степенной делитель, переставленные модули и пакет с несколькими необратимыми значениями.

\references
\end{document}
